\section{Часть 1. Анализ безопасности кода разработанного приложения}
\label{sec:part1}

\textbf{Краткое описание приложения.}
Разработано веб-приложение на языке Python с использованием фреймворка Flask. Все пользовательские записи живут только в адресном пространстве процесса и не сохраняются в файловой системе и СУБД. Интерфейс поддерживает создание, обновление и удаление двух классов данных: неконфиденциальных заметок (заголовок и текст хранятся открытым текстом) и конфиденциальных секретов (после ввода тело шифруется симметричной схемой Fernet, в модели остаётся только шифртекст). Ключ шифрования формируется при первом запуске, записывается с правами доступа только для владельца и не встраивается в исходный код. Сервер принимает соединения на локальном интерфейсе.

Интерфейс до и после ввода данных показан на рисунках~\ref{fig:ui} и~\ref{fig:filled}.

\begin{figure}[H]
	\centering
	\includegraphics[width=0.88\textwidth]{01_ui.png}
	\caption{Интерфейс приложения до ввода данных}
	\label{fig:ui}
\end{figure}

\begin{figure}[H]
	\centering
	\includegraphics[width=0.88\textwidth]{02_filled.png}
	\caption{Заметка и секрет после сохранения}
	\label{fig:filled}
\end{figure}

\textbf{Анализ защищённости кода.}
Безопасный код в смысле методички~\mdash  это код, который снижает риск нарушения конфиденциальности и доступности при невалидном вводе, сбоях и нагрузке. В разработанном приложении конфиденциальные и неконфиденциальные сущности разделены на уровне типов. Для секрета открытый текст принимается во временный буфер \texttt{bytearray}, сразу шифруется, буфер затирается нулями, после чего интерпретатор запрашивает сборку мусора. Обязательные поля и предельные длины проверяются на сервере; сообщения об ошибках не содержат введённого секрета. Журнал приложения не фиксирует тела конфиденциальных записей. На экране списка секретов отображается только заголовок и признак того, что содержимое зашифровано; расшифровка выполняется лишь по явному запросу пользователя и не записывается обратно в хранилище.

\begin{lstlisting}[caption={Шифрование секрета и затирание открытого текста},label={lst:wipe}]
buf = bytearray(plain.encode("utf-8"))
try:
    token = fernet.encrypt(bytes(buf))
finally:
    wipe_bytes(buf)
    gc.collect()
\end{lstlisting}

Ограничения языка остаются существенными. В CPython строки, поступившие из HTTP-формы, могут ещё некоторое время находиться в куче как копии буферов обработки запроса. Поэтому сразу после создания записи маркеры открытого текста иногда обнаруживаются в дампе, даже если в модели секрета уже лежит только шифртекст. Полностью исключить такие копии без нативных расширений нельзя; практическая цель реализации~\mdash  не удерживать открытый текст в долгоживущих структурах данных.

При аварийном завершении конфиденциальные данные не должны попадать на диск в виде журналов приложения. Дамп процесса, напротив, отражает содержимое ОЗУ и поэтому используется как инструмент анализа, а не как штатный канал хранения.

Ключ Fernet по умолчанию записывается в файл экземпляра с режимом 600. Это упрощает повторный запуск, но противоречит идеалу <<конфиденциальное только в ОЗУ>>: вместе с дампом процесса ключ с диска позволяет восстановить секреты. Для демонстрации без файла ключ можно передать переменной окружения; на учебном узле файл оставлен, каталог экземпляра закрыт для прочих uid.

\textbf{Анализ оперативной памяти и процессорного времени.}
Экспериментальный прогон выполнен 12.09.2026. По данным \texttt{/proc} и \texttt{ps} для работающего сервера: резидентный размер $\approx$ 37~МБ, виртуальный размер $\approx$ 128~МБ, один рабочий поток. Профилировщик запускался на стороне клиента (двадцать запросов состояния и записи): суммарное время около 0,24~с, основная доля~\mdash  сетевой стек клиента, а не шифрование на сервере. При учебном объёме данных криптография не является узким местом по процессорному времени, а рост памяти определяется прежде всего самим интерпретатором и каркасом веб-сервера.

\pagebreak
\textbf{Анализ дампов оперативной памяти.}
Дампы процесса снимались после создания, обновления и удаления записей утилитой \texttt{gcore}; поиск маркеров выполнялся по текстовым строкам дампа. В качестве неконфиденциального маркера использовалась заранее известная фраза заметки, в качестве конфиденциальных~\mdash  исходный и обновлённый тексты секрета. Сводка приведена в таблице~\ref{tab:dumps}.

\begin{table}[ht]
	\begin{tabularx}{\textwidth}{|l|c|c|c|}
		\hline
		\textbf{Сценарий} & Заметка & Исходный секрет & Обновлённый секрет \\
		\hline
		После создания & 1 & 3 & 0 \\
		\hline
		После обновления & 2 & 1 & 0 \\
		\hline
		После удаления & 1 & 1 & 0 \\
		\hline
	\end{tabularx}
	\caption{Число вхождений маркеров в дампах ОЗУ}
	\label{tab:dumps}
\end{table}

После создания открытый текст секрета ещё встречается в дампе (кратковременные копии запроса), но не хранится в модели данных. После обновления число вхождений исходного маркера падает, новый открытый текст в дампе не обнаружен. После удаления копия исходного маркера может остаться в буферах интерпретатора, хотя объект секрета из хранилища уже удалён. Неконфиденциальный маркер сохраняется дольше: для заметок это ожидаемо, поскольку они по постановке задачи хранятся открытым текстом. Таким образом, различие классов данных наблюдается не только в исходном коде, но и в фактическом содержимом оперативной памяти.
